%============================================================
%  Chapter 2 – Literature Review
%============================================================
\chapter{Literature Review and Related Work}

\section{Overview}

This chapter analyses existing freelance marketplace platforms and academic research relevant to the design of \projectname{}. The review identifies the strengths and limitations of current systems, establishing the research gap that this project aims to address.

\section{Existing Freelance Marketplace Platforms}

\subsection{Upwork}

Upwork \cite{upwork2024} is the world's largest freelance marketplace by revenue. It supports both hourly and fixed-price contracts and provides a comprehensive suite of tools including time-tracking, escrow payments, and a dispute resolution centre. Upwork's matching algorithm uses a proprietary ``Job Success Score'' (JSS) that combines client feedback, project outcomes, and account history.

\textbf{Limitations}: The platform charges up to 20\% freelancer commission on earnings below \$500, which many users consider excessive. The review system, while thorough, can be gamed by cancelling contracts before feedback is submitted. Non-technical clients still struggle with the proposal review process.

\subsection{Fiverr}

Fiverr \cite{fiverr2024} popularised the ``gig'' model, where freelancers create pre-packaged service offerings rather than bidding on client jobs. This model significantly lowers friction for clients who need a well-defined service.

\textbf{Limitations}: The gig model disadvantages custom, complex projects. Prices are anchored low, creating race-to-the-bottom dynamics. Review manipulation is a documented problem.

\subsection{Freelancer.com}

Freelancer.com \cite{freelancer2024} operates a competitive bidding model similar to \projectname{}. It features contests, milestone-based payments, and a large global user base.

\textbf{Limitations}: The platform is notorious for low-quality, low-cost proposals and a high percentage of spam bids. There is no AI-assisted scoping feature. Dispute resolution is slow and often unsatisfactory.

\subsection{Toptal}

Toptal \cite{toptal2024} takes the opposite approach --- an elite, invitation-only talent network where fewer than 3\% of applicants are accepted. This ensures quality but sacrifices accessibility.

\textbf{Limitations}: Not suitable for small projects or emerging-market freelancers. Opaque vetting process.

\section{Comparative Analysis}

\begin{table}[H]
\centering
\caption{Comparative Feature Analysis of Freelance Platforms}
\begin{tabularx}{\textwidth}{lXXXXX}
\toprule
\textbf{Feature} & \textbf{Upwork} & \textbf{Fiverr} & \textbf{Freelancer} & \textbf{Toptal} & \textbf{\projectname{}} \\
\midrule
AI Project Scoping    & No   & No   & No   & No   & \textbf{Yes} \\
Token Economy         & No   & No   & No   & No   & \textbf{Yes} \\
Blind Review System   & No   & No   & No   & No   & \textbf{Yes} \\
Milestone Contracts   & Yes  & Yes  & Yes  & Yes  & \textbf{Yes} \\
Dispute Resolution    & Yes  & Yes  & Yes  & Yes  & \textbf{Yes} \\
Google OAuth          & Yes  & Yes  & No   & No   & \textbf{Yes} \\
Admin Dashboard       & N/A  & N/A  & N/A  & N/A  & \textbf{Yes} \\
Role-Based Access     & Yes  & Yes  & Yes  & Yes  & \textbf{Yes} \\
Open-Source Stack     & No   & No   & No   & No   & \textbf{Yes} \\
\bottomrule
\end{tabularx}
\label{tab:platform_comparison}
\end{table}

\section{Academic Research}

\subsection{Token Economies in Online Platforms}

The concept of internal token economies has been extensively studied in the context of online communities and incentive design. Token-based systems have been shown to increase user engagement, reduce spam, and reward high-quality contributions \cite{designPatterns}. \projectname{} implements a \textit{SkillToken} economy where tokens are spent to submit proposals and awarded for registration and monthly activity, creating a self-regulating market for proposal quality.

\subsection{AI-Assisted Requirements Engineering}

The application of Large Language Models (LLMs) to requirements engineering is a growing research area. LLMs such as GPT-4 \cite{openai2024} have demonstrated the ability to transform informal natural-language descriptions into structured requirements artefacts. \projectname{} leverages this capability to help non-technical clients formulate project postings with appropriate category classifications, budget ranges, size estimates, and required skills.

\subsection{Blind Review Systems}

The ``blind review'' construct is borrowed from academic peer review, where reviewer identities are hidden to reduce bias. Applied to marketplace review systems, both-party-blind reviews --- where ratings are only revealed once both parties have submitted --- have been shown to produce more honest and balanced feedback.

\subsection{REST API Design Principles}

The architectural style of REST, originally defined by Fielding \cite{restApi}, underpins the entire backend API of \projectname{}. Key principles applied include stateless communication, resource-oriented URL design, standard HTTP method semantics, and JSON as the data interchange format.

\section{Identified Research Gaps}

Based on the above review, the following gaps justify the development of \projectname{}:

\begin{enumerate}
  \item No major freelance platform provides an \textbf{AI-assisted project scoping tool} integrated into the project posting workflow.
  \item No platform implements a \textbf{token-based proposal economy} to dissuade low-quality mass bidding.
  \item No platform implements a \textbf{fully blind mutual review system} where both reviews are hidden until simultaneous reveal.
  \item The academic community lacks a fully-documented, open-architecture reference implementation of a freelance marketplace for educational purposes.
\end{enumerate}

\projectname{} directly addresses all four of these gaps.
